\chapter{Analisi dei Requisiti e Casi d'uso}

Questo capitolo descrive il processo di progettazione del sistema, seguendo un approccio che va dall'esterno verso l'interno. Si parte dall'analisi di come gli utenti interagiscono con l'applicazione, definendo i requisiti e l'interfaccia grafica, per poi scendere progressivamente nei dettagli tecnici dell'architettura software e della struttura del database.

\section{Use Case Diagram}
Il diagramma dei casi d'uso offre una visione d'insieme delle funzionalità del sistema. Serve a mappare visivamente quali azioni possono compiere i tre attori principali (Driver, Station Operator ed Energy Manager) all'interno della piattaforma.

\begin{figure}[H]
    \centering
    \includegraphics[width=\textwidth]{UseCaseDiagramNew.pdf.drawio.pdf}
    \caption{Diagramma dei casi d'uso}
    \label{fig:mio_schema}
\end{figure}

\section{Use Case Templates}
Per evitare ambiguità, i casi d'uso più importanti sono stati documentati nel dettaglio attraverso dei template strutturati. Questi schemi descrivono passo dopo passo il flusso degli eventi, le precondizioni necessarie e i possibili scenari di errore.
\begingroup
\renewcommand{\arraystretch}{1.4}
\begin{xltabular}{\textwidth}{| >{\bfseries}l | X |}

% --- INTESTAZIONE PRIMA PAGINA ---
\caption{Caso d'Uso: Ricarica Veicolo} \label{tab:uc_ricarica} \\
\hline
\rowcolor{gray!15} Campo & \textbf{Contenuto} \\ \hline
\endfirsthead

% --- INTESTAZIONE PAGINE SUCCESSIVE ---
\multicolumn{2}{c}{{\tablename\ \thetable{} -- Continua dalla pagina precedente}} \\
\hline
\rowcolor{gray!15} Campo & \textbf{Contenuto} \\ \hline
\endhead

% --- PIÈ DI PAGINA (se si spezza) ---
\hline
\multicolumn{2}{r}{\textit{Continua nella pagina successiva...}} \\
\endfoot

% --- PIÈ DI PAGINA FINALE ---
\endlastfoot

% --- CONTENUTO ---
Nome & Ricarica Veicolo (Charge Vehicle) \\ \hline
Livello & Obiettivo Utente (User Goal) \\ \hline
Descrizione & Il Conducente avvia una sessione di ricarica su una stazione selezionata, scegliendo una specifica strategia di ricarica (Fast o Eco). La sessione viene eseguita in tempo reale: il Grid Monitor la fa avanzare e l'interfaccia utente (Live Session UI) ne osserva lo stato tramite polling. \\ \hline
Pre-condizioni & \parbox[t]{\linewidth}{%
Il Conducente ha effettuato l'accesso.\\
Il saldo del Portafoglio (Wallet) del Conducente copre almeno il costo di 5 kWh presso la stazione selezionata, calcolato per il piano del Conducente (\texttt{station.priceFor(5.0, driver)}) e la strategia selezionata.\\
La stazione selezionata è \texttt{ACTIVE} (oppure \texttt{RESERVED} dal Conducente) ed è compatibile con il connettore del veicolo.
} \\ \hline
Post-condizioni & La sessione è \texttt{COMPLETED}, il costo (detratto in modo incrementale per ogni tick durante la ricarica) è interamente addebitato sul Portafoglio, viene registrata una Transazione di tipo \texttt{CHARGE} e la stazione torna allo stato \texttt{ACTIVE}. \\ \hline
Attori & Conducente, Grid Monitor (Sistema) \\ \hline
Unit Test & \texttt{SessionServiceTest} \\ \hline
Flusso Principale & 
\begin{enumerate}[leftmargin=*, nosep]
    \item Il Conducente seleziona una stazione dalla lista delle stazioni disponibili e compatibili più vicine (\texttt{ACTIVE}, o \texttt{RESERVED} da lui stesso).
    \item Il Conducente apre la Live Session per quella stazione.
    \item Il Conducente sceglie una strategia di ricarica: Fast Charge o Eco Charge.
    \item Il Conducente inserisce la percentuale di batteria attuale del veicolo.
    \item Il Conducente preme Start; il sistema valida il saldo del Portafoglio e la disponibilità della stazione, imposta la stazione su \texttt{BUSY} e apre la sessione.
    \item Il Grid Monitor, ogni 5s, aggiorna batteria e kWh erogati, accumula il calore sul trasformatore di alimentazione e deduce i fondi in base alla strategia (lo sconto del piano si applica solo alla quota tariffaria della piattaforma).
    \item La Live Session UI interroga (poll) lo stato della sessione e riflette in tempo reale batteria, kWh, costo attuale e Portafoglio.
    \item Il Conducente preme Stop (oppure la sessione si auto-completa quando la batteria raggiunge il 100\%).
    \item Il sistema calcola i kWh totali erogati, il costo finale e registra la Transazione.
    \item Il sistema reimposta lo stato della stazione su \texttt{ACTIVE} e mostra il popup di valutazione (rating).
\end{enumerate} \\ \hline
Flussi Alternativi & 
\textbf{1. Stazioni non disponibili nascoste:} \newline
1.1. Le stazioni \texttt{BUSY}, \texttt{OVERLOADED}, incompatibili o \texttt{RESERVED} da altri non sono mostrate in lista, impedendo al Conducente di selezionare una stazione non disponibile. \newline\newline
\textbf{5. Saldo insufficiente:} \newline
5.1. Se il saldo è inferiore al costo di 5 kWh, il sistema nega l'avvio e mostra: "Credito insufficiente per avviare la sessione." (\texttt{InsufficientBalanceException}). \newline\newline
\textbf{6-7. Allarme Termico (Observer attivato):} \newline
6.1. Se la temperatura del trasformatore supera i 90°C, il \texttt{LoadBalancerService} (Observer) imposta tutte le stazioni su quel trasformatore su \texttt{OVERLOADED} e chiude forzatamente le relative sessioni attive. \newline
6.2. La sessione diventa \texttt{INTERRUPTED\_THERMAL} e la UI mostra: "Arresto di emergenza: surriscaldamento del sistema". \newline
6.3. Il sistema elabora un rimborso automatico per i kWh pagati ma non erogati. \\ \hline

\end{xltabular} 
\endgroup

\begingroup
\renewcommand{\arraystretch}{1.4}
\begin{xltabular}{\textwidth}{| >{\bfseries}l | X |}

% --- INTESTAZIONE PRIMA PAGINA ---
\caption{Caso d'Uso: Prenotazione a Breve Termine} \label{tab:uc_prenotazione} \\
\hline
\rowcolor{gray!15} Campo & \textbf{Contenuto} \\ \hline
\endfirsthead

% --- INTESTAZIONE PAGINE SUCCESSIVE ---
\multicolumn{2}{c}{{\tablename\ \thetable{} -- Continua dalla pagina precedente}} \\
\hline
\rowcolor{gray!15} Campo & \textbf{Contenuto} \\ \hline
\endhead

% --- PIÈ DI PAGINA (se si spezza) ---
\hline
\multicolumn{2}{r}{\textit{Continua nella pagina successiva...}} \\
\endfoot

% --- PIÈ DI PAGINA FINALE ---
\endlastfoot

% --- CONTENUTO ---
Nome & Prenotazione a Breve Termine \\ \hline
Livello & Obiettivo Utente \\ \hline
Descrizione & Il Conducente riserva per 15 minuti una colonnina \texttt{ACTIVE}, per assicurarsene la disponibilità al proprio arrivo. \\ \hline
Pre-condizioni & \parbox[t]{\linewidth}{%
Il Conducente ha effettuato l'accesso.\\
Non ha già un'altra prenotazione attiva.\\
La colonnina scelta è \texttt{ACTIVE} ed è compatibile con il connettore del veicolo.
} \\ \hline
Post-condizioni & La colonnina è in stato \texttt{RESERVED}, il campo \texttt{reserved\_by\_id} punta al Conducente e \texttt{expiration\_timestamp} è impostato a 15 minuti dall'istante corrente. La colonnina resta visibile al solo Conducente che l'ha prenotata. \\ \hline
Attori & Conducente \\ \hline
Unit Test & \texttt{StationServiceTest} \\ \hline
Flusso Principale &
\begin{enumerate}[leftmargin=*, nosep]
    \item Il Conducente seleziona una colonnina \texttt{ACTIVE} dall'elenco di quelle disponibili.
    \item Conferma di volerla prenotare.
    \item Il sistema tenta l'acquisizione con un \texttt{UPDATE} condizionale, che modifica la riga solo se la colonnina risulta ancora \texttt{ACTIVE}.
    \item Il sistema porta la colonnina a \texttt{RESERVED} e ne registra il Conducente come titolare.
    \item Il sistema imposta \texttt{expiration\_timestamp} a 15 minuti dall'istante corrente.
    \item Il sistema conferma la prenotazione mostrando l'orario esatto di scadenza.
\end{enumerate} \\ \hline
Flussi Alternativi &
\textbf{3. Colonnina non più disponibile:} \newline
3.1. Se nel frattempo un altro Conducente l'ha prenotata o occupata, l'\texttt{UPDATE} non aggiorna alcuna riga e il sistema nega la prenotazione. \newline
3.2. La transazione viene annullata con un \texttt{rollback} e l'oggetto in memoria non viene modificato, così che memoria e database non divergano. \newline
3.3. Il sistema mostra il messaggio \textit{"La colonnina non è più disponibile"} e si torna al passo 1. \newline\newline
\textbf{Nota sulla concorrenza:} \newline
\textbullet\ È il database stesso a impedire la doppia prenotazione: la condizione sullo stato \texttt{ACTIVE} è parte dell'\texttt{UPDATE}, quindi fra due Conducenti simultanei ne può vincere uno solo, senza alcun blocco applicativo. \newline\newline
\textbf{Nota sulla scadenza:} \newline
\textbullet\ Se i 15 minuti trascorrono senza che venga avviata una sessione, il \texttt{GridMonitor} riporta la colonnina ad \texttt{ACTIVE} (Tab.~\ref{tab:uc_update_state}). \\ \hline

\end{xltabular}
 
\endgroup

\section{Mockups}
In questa fase sono stati realizzati i mockup dell'interfaccia grafica per verificare la User Experience e garantire
che tutte le funzionalità descritte nei casi d'uso siano accessibili in modo immediato ed intuitivo.

\section{Page Navigation Diagram}
Il diagramma di navigazione illustra in che modo le diverse schermate dell'applicazione sono collegate tra loro. Mostra il percorso che l'utente compie muovendosi all'interno del sistema in base alle scelte effettuate o ai pulsanti premuti.

\begin{figure}[H]
    \centering
    \includegraphics[width=\textwidth]{NavigationPageDiagram.drawio2.pdf}
    \caption{Diagramma di navigazione tra le pagine dell'applicazione.}
    \label{fig:navigation_schema}
\end{figure}

\section{Package Diagram}
Passando all'architettura interna del software, il diagramma dei pacchetti mostra come il codice è stato diviso in moduli logici (come ad esempio la logica di business, l'interfaccia e l'accesso ai dati). L'obiettivo è mantenere il progetto ordinato e limitare le dipendenze tra i vari componenti.

\section{Class Diagram}
Il sistema è stato progettato con un'architettura a strati per mantenere alta coesione e basso accoppiamento (\textit{high coesion and low coupling}). Nel seguente paragrafo sono approfondite le classi principali:

\subsection{Domain Model}

Questo diagramma rappresenta il Domain Model, cioè le entità principali su cui si basa il sistema: utenti, colonnine, trasformatori, sessioni di ricarica e gli oggetti ad esse collegati. Una descrizione più dettagliata sul Domain Model è presente nel paragrafo \ref{sec: Domain Model}

\begin{figure}[H]
    \centering
    \includegraphics[width=1.10\textwidth]{domain_model.png}
    \caption{Domain Model, diagramma di classe}
    \label{fig:domain_model}
\end{figure}

\subsection{Business Layer}

Il seguente diagramma mostra il Business Layer, il livello che contiene la logica applicativa del sistema. Qui si trovano i servizi che coordinano le entità del Domain Model, le strategie di ricarica e i componenti che gestiscono la simulazione della rete. Il livello è approfondito nel paragrafo \ref{sec: Business}.

\begin{figure}[H]
    \centering
    \includegraphics[width=1.10\textwidth]{buiness.png}
    \caption{Business Layer, diagramma di classe}
    \label{fig:business_layer}
\end{figure}

\subsection{DAO}

Di seguito è mostrato il diagramma del livello DAO, responsabile della comunicazione con il database. Il suo compito è tradurre le entità del Domain Model in dati salvabili e recuperabili, garantendone quindi la persistenza

\begin{figure}[H]
    \centering
    \includegraphics[width=1.10\textwidth]{dao.png}
    \caption{DAO, diagramma di classe}
    \label{fig:DAO}
\end{figure}

\subsection{Presentation Layer}

Il seguente diagramma rappresenta il presentation layer, cioè l'interfaccia grafica dell'applicazione realizzata in JavaFX. Comprende i controller delle varie schermate, che si occupano di mostrare i dati e raccogliere le azioni dell'utente, appoggiandosi ai servizi del livello di business.

\begin{figure}[H]
    \centering
    \includegraphics[width=1.10\textwidth]{presentation.png}
    \caption{Presentation Layer, diagramma di classe}
    \label{fig:presentation_layer}
\end{figure}

\section{Sequence Diagram}

Nel seguente paragrafo sono mostrati i diagrammi di sequenza di alcuni dei flussi più significativi del sistema, che descrivono come le varie classi interagiscono tra loro nel tempo per portare a termine un'operazione.

\subsection{Ciclo di vita del GridMonitor
}
\begin{figure}[H]
    \centering
    \includegraphics[width=1.10\textwidth]{GridMonitor_Lifecycle.png}
    \caption{Diagramma di sequenza del ciclo di vita del \texttt{GridMonitor}.}
    \label{fig:seq_gridmonitor}
\end{figure}

Il \texttt{GridMonitor} (Fig.~\ref{fig:seq_gridmonitor}) è il thread daemon che scandisce l'intera simulazione della rete elettrica con un tick fisso di 5 secondi, ed è di fatto l'unico componente attivo di propria iniziativa nel sistema: tutti gli altri servizi vengono invocati su richiesta. Ad ogni ciclo il monitor libera prima le prenotazioni scadute tramite \texttt{StationService.expireHolds()}, poi recupera da \texttt{SessionService} l'elenco delle sessioni attive e, per ciascuna, seleziona la \texttt{ChargingStrategy} corrispondente per calcolare kWh erogati e costo del tick, aggiornando in un'unica transazione sia il wallet del Driver sia lo stato della sessione. Lo stesso tick genera un incremento di calore che viene accumulato per trasformatore e applicato a fine ciclo a ciascun \texttt{PowerTransformer}: se la temperatura supera i 90°C viene notificato l'evento \texttt{THERMAL\_ALERT} agli observer registrati, dando origine al flusso descritto in Fig.~\ref{fig:seq_thermal_alert}. Infine il monitor invoca il controllo di qualità del \texttt{RatingService}. È importante osservare che la UI non partecipa a questo ciclo come observer: lo stato delle sessioni viene mostrato tramite un polling separato sul database, in modo tale che i due flussi restino disaccoppiati e l'unico punto di sincronizzazione sia la persistenza.

\subsection{Approvazione di una colonnina}

\begin{figure}[H]
    \centering
    \includegraphics[width=0.9\textwidth]{Station_Approval.png}
    \caption{Diagramma di sequenza dell'approvazione di una colonnina da parte dell'Energy Manager.}
    \label{fig:seq_station_approval}
\end{figure}

Il flusso di approvazione (Fig.~\ref{fig:seq_station_approval}) regola l'ingresso in esercizio di una nuova colonnina proposta da uno Station Operator. L'Energy Manager apre la coda delle richieste in stato \texttt{PENDING\_VERIFICATION}, popolata da \texttt{StationService.getPendingStations()}, e per ciascuna voce può approvare o rifiutare. In caso di approvazione il controller richiede la tariffa di piattaforma da applicare in €/kWh e delega la transizione alla \texttt{ValidationService}, che apre esplicitamente una connessione con \texttt{setAutoCommit(false)} prima di invocare \texttt{activate()} sul dominio: è la \texttt{ChargingStation} stessa a validare che la transizione sia legittima solo a partire da \texttt{PENDING\_VERIFICATION}, sollevando un'eccezione altrimenti, mentre il servizio si occupa unicamente di orchestrare la persistenza e il commit o il rollback della transazione in caso di errore. La stessa separazione si ritrova nel percorso di rifiuto, dove il dominio applica la transizione a \texttt{REJECTED} dopo l'acquisizione della motivazione da parte del Manager. In entrambi i casi la responsabilità di stabilire se un cambio di stato è ammesso resta interamente nel Domain Model , mentre il Business Layer si limita a decidere quando invocarla e a garantirne l'atomicità della persistenza.

\subsection{Gestione dell'evento di temperatura troppo alta (Thermal Alert)}
\begin{figure}[H]
    \centering
    \includegraphics[width=0.9\textwidth]{Thermal_Alert.png}
    \caption{Diagramma di sequenza della gestione di una temperatura troppo alta (Thermal Alert).}
    \label{fig:seq_thermal_alert}
\end{figure}

Il diagramma di Thermal Alert (Fig.~\ref{fig:seq_thermal_alert}) mostra cosa succede quando un trasformatore si surriscalda. Durante il tick del \texttt{GridMonitor}, il \texttt{PowerTransformer} riceve il nuovo carico di calore e aggiorna la propria temperatura interna. Se questa supera i 90°C, il trasformatore avvisa il \texttt{LoadBalancerService}, che è l'unico componente incaricato di reagire all'emergenza: il trasformatore stesso non sa cosa succederà dopo, si limita a segnalare il problema.

A quel punto il \texttt{LoadBalancerService} interviene in due passaggi separati. Nel primo, blocca subito tutte le colonnine collegate a quel trasformatore mettendole in stato \texttt{OVERLOADED}, per evitare che qualcuno continui a ricaricare su un ramo elettrico in sovraccarico: questo blocco viene fatto tutto insieme, colonnina per colonnina, come un'unica operazione che riesce o fallisce nel suo complesso. Nel secondo passaggio, il sistema chiude una per una tutte le sessioni di ricarica che erano attive su quelle colonnine, rimborsando ogni Driver coinvolto tramite \texttt{SessionService.forceClose()}. Qui però ogni chiusura viene trattata singolarmente: se il rimborso di una sessione dovesse fallire, le altre sessioni verrebbero comunque chiuse e rimborsate correttamente, e soprattutto le colonnine resterebbero comunque bloccate in sicurezza, indipendentemente da come va il rimborso.

\section{Database}

La gestione dei dati di ChargeNet è affidata a un database relazionale PostgreSQL. Inizialmente, per velocizzare il lavoro, si è scelto di lavorare inizialmente con H2. Questo ha permesso di scrivere query e testare la parte applicativa fin da subito, per poi occuparsi successivamente dell'installazione e configurazione del database reale, passaggio che è stato comunque rapido visto che non è stato necessario modificare la sintassi SQL (H2 è stato usato in modalità \texttt{MODE=PostgreSQL})

\subsection{Schema ER}

\begin{figure}[H]
    \centering
    \includegraphics[width=0.95\textwidth]{ChargeNet_ER.drawio.pdf}
    \caption{Schema Entità-Relazione del database ChargeNet.}
    \label{fig:er_chargenet}
\end{figure}

Lo schema (Fig.~\ref{fig:er_chargenet}) riflette la stessa gerarchia già introdotta nel Domain Model: \texttt{Users} è un'entità generica che rappresenta ciò che ogni utente ha in comune (identità, credenziali, ruolo), che viene specializzata in tre sottotipi,\textbf{Driver}, \textbf{StationOperator} ed \textbf{EnergyManager}, ciascuno con i propri attributi specifici. Un utente appartiene esattamente a uno dei tre sottotipi, mai a più di uno contemporaneamente, infatti la colonna \texttt{role} può avere un solo valore.

Le relazioni principali partono dallo specifico sottotipo e non genericamente da \texttt{Users}:

\begin{itemize}
    \item Un \textbf{PowerTransformer} alimenta più \textbf{ChargingStation}.
    \item Uno \textbf{StationOperator} possiede più \textbf{ChargingStation}.
    \item Un \textbf{Driver} può prenotare temporaneamente una colonnina, avviare più \textbf{ChargingSession}, generare più \textbf{Transaction} e lasciare più \textbf{Rating}.
    \item Una \textbf{ChargingStation} ospita più \textbf{ChargingSession} e può ricevere più \textbf{Rating} o generare più \textbf{RatingAlert}.
    \item Ogni sessione può generare al più una \textbf{Rating}.
    \item \textbf{EnergyManager} non compare in nessuna relazione dello schema: il suo ruolo (approvare colonnine, gestire gli alert di qualità) non richiede di possedere o generare entità proprie, ma di agire su entità già esistenti.
\end{itemize}

\subsection{Modello Relazionale}

La traduzione delle entità semplici segue le regole standard: ogni entità diventa una tabella, ogni relazione (1,n) diventa una chiave esterna sul lato "molti". La generalizzazione invece è stata tradotta nel seguente modo: un'unica tabella con tutte le colonne di tutti i sottotipi, dove le colonne che non riguardano una riga restano \texttt{NULL}, e una colonna discriminatore indica il sottotipo reale. 

\begin{quote}
\ttfamily
Users(\uline{id}, name, email, password, role, latitude, longitude, wallet\_balance,\\
\hspace*{3.2em}connector\_type, subscription\_plan, battery\_capacity, total\_earnings)

\vspace{0.8em}
PowerTransformers(\uline{id}, name, temperature, load\_percent)

\vspace{0.8em}
ChargingStations(\uline{id}, \dashuline{operator\_id}, \dashuline{transformer\_id}, name, address,\\
\hspace*{3.2em}latitude, longitude, connector\_type, power\_kw, is\_solar\_powered,\\
\hspace*{3.2em}tariff\_operator, tariff\_platform, average\_rating, total\_ratings,\\
\hspace*{3.2em}status, \dashuline{reserved\_by\_id}$^{\ast}$, expiration\_timestamp)

\vspace{0.8em}
ChargingSessions(\uline{id}, \dashuline{driver\_id}, \dashuline{station\_id}, strategy\_used,\\
\hspace*{3.2em}battery\_start, battery\_current, kwh\_delivered, cost\_total, status,\\
\hspace*{3.2em}opened\_at, closed\_at$^{\ast}$)

\vspace{0.8em}
Transactions(\uline{id}, \dashuline{driver\_id}, type, amount, kwh, description, created\_at)

\vspace{0.8em}
Ratings(\uline{id}, \dashuline{driver\_id}, \dashuline{station\_id}, \dashuline{session\_id}, stars,\\
\hspace*{3.2em}comment, created\_at) \quad \textsc{unique}(driver\_id, session\_id)

\vspace{0.8em}
RatingAlerts(\uline{id}, \dashuline{station\_id}, avg\_at\_creation, status, manager\_note$^{\ast}$, created\_at)
\end{quote}

\vspace{0.3em}
\noindent dove:
\begin{quote}
\ttfamily
operator\_id, reserved\_by\_id \ \ si riferisce a \ \ Users(id)\\
transformer\_id \ \ si riferisce a \ \ PowerTransformers(id)\\
driver\_id \ \ si riferisce a \ \ Users(id)\\
station\_id \ \ si riferisce a \ \ ChargingStations(id)\\
session\_id \ \ si riferisce a \ \ ChargingSessions(id)
\end{quote}

\vspace{1em}
\noindent\textbf{Legenda della notazione}
\begin{itemize}
    \item[] \uline{attributo} \ — \ Chiave Primaria (PK) della tabella.
    \item[] \dashuline{attributo} \ — \ Chiave Esterna (FK) verso un'altra tabella.
    \item[] $^{\ast}$ \ — \ Attributo opzionale (\texttt{NULL} ammesso); tutti gli altri attributi sono obbligatori (\texttt{NOT NULL}).
    \item[] \textsc{unique}(\dots) \ — \ Vincolo di unicità su una combinazione di più colonne.
\end{itemize}

La colonna \texttt{role} in \texttt{Users} dice se la riga rappresenta un Driver, uno StationOperator o un EnergyManager. Colonne come \texttt{wallet\_balance} (Driver) o \texttt{total\_earnings} (StationOperator) esistono per tutte le righe della tabella, ma restano \texttt{NULL} per i sottotipi a cui non appartengono. Per esempio ovviamente non ha senso che \texttt{EnergyManager} abbia un valore diverso da \texttt{null} nell'attributo \texttt{subscription\_plan}. È lo stesso meccanismo già visto nel paragrafo ORM riguardo il metodo \texttt{extractUserFromResultSet()}, che legge la colonna \texttt{role} e ricostruisce l'oggetto Java del sottotipo corretto.


Per quanto riguarda invece le due FK verso \texttt{Users} in \texttt{ChargingStations} (\texttt{operator\_id} e \texttt{reserved\_by\_id}) modellano due relazioni indipendenti — chi possiede la colonnina e chi la sta eventualmente prenotando, ed è proprio per questo che solo la seconda può essere nulla.

Il vincolo \texttt{UNIQUE (driver\_id, session\_id)} sulla tabella \texttt{Ratings} fa si che anche a livello di database valga la regola che un Driver può lasciare una sola valutazione per ciascuna sessione.

Infine, diversi attributi numerici sono protetti da vincoli \texttt{CHECK} direttamente nello schema (es. \texttt{stars} tra 1 e 5, \texttt{wallet\_balance} sempre $\geq 0$): questi vincoli sono una seconda difesa rispetto alla validazione già presente nel Domain Model, indipendente da eventuali bug nel codice Java.

E' presente un limite, conseguenza della generalizzazione utilizzata: la foreign key \texttt{operator\_id} garantisce solo che l'id referenziato esista in \texttt{Users}, non che quella riga appartenga effettivamente a uno StationOperator, lo stesso vale per \texttt{driver\_id} rispetto ai Driver. Questo vincolo di coerenza tra ruolo e utilizzo non è esprimibile a livello di schema relazionale standard, tuttavia è comunque garantito nel codice Java. Infatti il metodo \texttt{StationService.registerStation()} accetta come parametro un \texttt{StationOperator}, non un generico \texttt{User}: un Driver non potrebbe essere passato a quel metodo nemmeno volendo, perché il codice non compilerebbe. La sicurezza che manca allo schema relazionale è quindi comunque garantita ma non appunto da quest'ultimo.